\documentclass[12pt,a4paper]{report}

% -----------------------------------------------------------------------------
% PACKAGES & SETUP (TCET 2026-27 GUIDELINES)
% -----------------------------------------------------------------------------
\usepackage[a4paper, left=1.25in, right=1.0in, top=1.0in, bottom=1.0in]{geometry}
\usepackage{setspace}
\onehalfspacing % 1.5 Line Spacing

% Fonts: Times New Roman Equivalent
\usepackage{mathptmx}
\usepackage[T1]{fontenc}
\usepackage[utf8]{inputenc}

% Math & Symbols
\usepackage{amsmath,amssymb,amsfonts}

% Graphics & Tables
\usepackage{graphicx}
\usepackage{booktabs}
\usepackage{tabularx}
\usepackage{longtable}
\usepackage{array}
\usepackage{float}
\usepackage{multirow}

% Headers, Footers & Formatting
\usepackage{fancyhdr}
\usepackage{titlesec}
\usepackage{cite}
\usepackage{listings}
\usepackage{xcolor}
\usepackage{hyperref}

\hypersetup{
    colorlinks=true,
    linkcolor=black,
    citecolor=black,
    urlcolor=blue
}

% Code Listing Style
\lstset{
    basicstyle=\ttfamily\footnotesize,
    breaklines=true,
    frame=single,
    backgroundcolor=\color{gray!6},
    numbers=left,
    numberstyle=\tiny\color{gray},
    keywordstyle=\color{blue}\bfseries,
    commentstyle=\color{green!50!black}\itshape,
    stringstyle=\color{purple},
    tabsize=2,
    showstringspaces=false
}

% Heading Styles (TCET Specifications)
\titleformat{\chapter}[display]
  {\normalfont\fontsize{18}{22}\bfseries\centering}{\chaptertitlename\ \thechapter}{14pt}{\fontsize{18}{22}\bfseries\MakeUppercase}
\titlespacing*{\chapter}{0pt}{-20pt}{20pt}

\titleformat{\section}
  {\normalfont\fontsize{16}{20}\bfseries}{\thesection}{1em}{}
\titlespacing*{\section}{0pt}{14pt}{6pt}

\titleformat{\subsection}
  {\normalfont\fontsize{14}{18}\bfseries}{\thesubsection}{1em}{}
\titlespacing*{\subsection}{0pt}{10pt}{4pt}

\titleformat{\subsubsection}
  {\normalfont\fontsize{12}{16}\bfseries}{\thesubsubsection}{1em}{}
\titlespacing*{\subsubsection}{0pt}{6pt}{2pt}

% Paragraph settings
\setlength{\parindent}{0pt}
\setlength{\parskip}{6pt}

% -----------------------------------------------------------------------------
% DOCUMENT BODY
% -----------------------------------------------------------------------------
\begin{document}

% -----------------------------------------------------------------------------
% 1. TITLE PAGE
% -----------------------------------------------------------------------------
\begin{titlepage}
\begin{center}

% -----------------------------------------------------------------------------
% TCET INSTITUTIONAL LOGO RESERVED SPACE
% To insert your college logo, upload tcet_logo.png to Overleaf and uncomment:
% \includegraphics[width=2.5cm]{tcet_logo.png}
\vspace*{2.0cm} % Reserved blank space for college logo
% -----------------------------------------------------------------------------

\vspace{0.8cm}
{\fontsize{14}{18}\selectfont \textbf{Project Report (Part I)}\par}
{\fontsize{14}{18}\selectfont on\par}
\vspace{0.8cm}

{\fontsize{22}{26}\selectfont \textbf{SAMVADA: Real-Time Multilingual Speech-to-Speech Translation Platform for Web-Based Video Conferencing}\par}

\vspace{1.0cm}
{\fontsize{12}{16}\selectfont \textit{Submitted in partial fulfillment for the award of the degree of}\par}
\vspace{0.2cm}
{\fontsize{16}{20}\selectfont \textbf{BACHELOR OF ENGINEERING}\par}
\vspace{0.1cm}
{\fontsize{12}{16}\selectfont \textit{In}\par}
\vspace{0.1cm}
{\fontsize{16}{20}\selectfont \textbf{COMPUTER ENGINEERING}\par}

\vspace{1.0cm}
{\fontsize{12}{16}\selectfont \textbf{Submitted by}\par}
\vspace{0.2cm}
{\fontsize{12}{16}\selectfont
Name of Group Member 1 (Roll No: \rule{3cm}{0.4pt})\\
Name of Group Member 2 (Roll No: \rule{3cm}{0.4pt})\\
Name of Group Member 3 (Roll No: \rule{3cm}{0.4pt})\par}

\vspace{1.0cm}
{\fontsize{12}{16}\selectfont \textbf{\textit{Under the Guidance of}}\par}
\vspace{0.2cm}
{\fontsize{14}{18}\selectfont \textbf{[Guide Name]}\par}
{\fontsize{14}{18}\selectfont \textbf{[Designation]}\par}

\vfill

{\fontsize{14}{18}\selectfont \textbf{Department of Computer Engineering}\\
\textbf{Thakur College of Engineering and Technology, Mumbai}\\
\textit{(An Autonomous College affiliated to University of Mumbai)}\\
\textbf{(Academic Year 2026-27)}\par}

\end{center}
\end{titlepage}

% -----------------------------------------------------------------------------
% 2. CERTIFICATE
% -----------------------------------------------------------------------------
\newpage
\thispagestyle{empty}
\begin{center}

% Reserved space for college logo on Certificate
\vspace*{1.5cm} % \includegraphics[width=2.2cm]{tcet_logo.png}

\vspace{0.5cm}
{\fontsize{16}{20}\selectfont \textbf{CERTIFICATE}\par}
\end{center}

\vspace{0.8cm}
\noindent This is to certify that the project entitled \textbf{``SAMVADA: Real-Time Multilingual Speech-to-Speech Translation Platform for Web-Based Video Conferencing''} is a bonafide work of \textbf{Name of Group Member 1 (Roll No: \rule{2cm}{0.4pt}), Name of Group Member 2 (Roll No: \rule{2cm}{0.4pt}), and Name of Group Member 3 (Roll No: \rule{2cm}{0.4pt})} submitted to the \textbf{Thakur College of Engineering and Technology, Mumbai} (An Autonomous College affiliated to University of Mumbai) in partial fulfillment of the requirement for the Project-I for award of the degree of \textbf{``Bachelor of Engineering''} in \textbf{``Computer Engineering''}.

\vspace{2.5cm}

\noindent
\begin{tabular}{p{7.5cm} p{7.5cm}}
Signature with Date: \hrulefill & Signature with Date: \hrulefill \\[1.2cm]
\textbf{Name of Guide:} [Guide Name] & \textbf{Name of HOD:} Dr. Vaishali Kaiche \\
\textbf{Designation:} [Designation] & \textbf{Department:} Computer Engineering \\
\end{tabular}

% -----------------------------------------------------------------------------
% 3. PROJECT APPROVAL CERTIFICATE
% -----------------------------------------------------------------------------
\newpage
\thispagestyle{empty}
\begin{center}

% Reserved space for college logo on Approval Certificate
\vspace*{1.5cm} % \includegraphics[width=2.2cm]{tcet_logo.png}

\vspace{0.5cm}
{\fontsize{16}{20}\selectfont \textbf{PROJECT APPROVAL CERTIFICATE}\par}
\end{center}

\vspace{0.8cm}
\noindent This project report entitled \textbf{``SAMVADA: Real-Time Multilingual Speech-to-Speech Translation Platform for Web-Based Video Conferencing''} by \textbf{Name of Student 1 (Roll No: \rule{2cm}{0.4pt}), Name of Student 2 (Roll No: \rule{2cm}{0.4pt}), and Name of Student 3 (Roll No: \rule{2cm}{0.4pt})} is approved for the degree of \textbf{``Bachelor of Engineering''} in \textbf{``Computer Engineering''}.

\vspace{2.5cm}

\noindent
\begin{tabular}{p{7.5cm} p{7.5cm}}
\textbf{Internal Examiner:} & \textbf{External Examiner:} \\[1.2cm]
Signature: \hrulefill & Signature: \hrulefill \\
Name: \hrulefill & Name: \hrulefill \\
Date: \hrulefill & Date: \hrulefill \\
Place: Mumbai & Place: Mumbai \\
\end{tabular}

% -----------------------------------------------------------------------------
% 4. ACKNOWLEDGEMENT
% -----------------------------------------------------------------------------
\newpage
\thispagestyle{empty}
\begin{center}
{\fontsize{16}{20}\selectfont \textbf{ACKNOWLEDGEMENT}\par}
\end{center}

\vspace{0.8cm}
\noindent It would be unfair if I do not acknowledge the help and support given by Professors, students, friends, and colleagues.

\noindent We sincerely thank our guide \textbf{[Guide Name]} for his/her invaluable guidance, technical critique, constant encouragement, and steadfast support throughout the realization of this project. We also thank the project coordinators for arranging the computational facilities and testing environment necessary to carry out this research.

\noindent We express our deepest gratitude to the Head of Department, \textbf{Dr. Vaishali Kaiche}, the Principal, \textbf{Dr. B. K. Mishra}, and the Management of Thakur College of Engineering and Technology for fostering an inspiring research ecosystem.

\vspace{2.5cm}

\noindent
\begin{flushright}
1. \rule{6cm}{0.4pt}\\
2. \rule{6cm}{0.4pt}\\
3. \rule{6cm}{0.4pt}\\
\textit{(Name of Student and Roll No.)}
\end{flushright}

% -----------------------------------------------------------------------------
% 5. PLAGIARISM REPORT PLACEHOLDER
% -----------------------------------------------------------------------------
\newpage
\thispagestyle{empty}
\begin{center}
{\fontsize{16}{20}\selectfont \textbf{Blue Book Plagiarism Report}\par}
{\fontsize{14}{18}\selectfont \textbf{(From department Turnitin account only)}\par}
\end{center}

\vspace{4.0cm}
\begin{center}
\textit{(Departmental Turnitin Plagiarism Clearance Certificate to be affixed here)}
\end{center}

% -----------------------------------------------------------------------------
% 6. FRONT MATTER LISTINGS
% -----------------------------------------------------------------------------
\newpage
\pagenumbering{roman}
\setcounter{page}{1}

\tableofcontents
\newpage
\listoffigures
\newpage
\listoftables

% -----------------------------------------------------------------------------
% 7. ABSTRACT
% -----------------------------------------------------------------------------
\newpage
\begin{center}
{\fontsize{16}{20}\selectfont \textbf{ABSTRACT}\par}
\end{center}

\vspace{0.5cm}
\noindent In an increasingly globalized world, cross-lingual verbal communication remains a formidable bottleneck in enterprise collaboration, academic research, healthcare delivery, and international crisis management. Contemporary video conferencing platforms (such as Zoom, Google Meet, and Microsoft Teams) either force participants into an artificial \textit{lingua franca} (predominantly English), mandate prohibitive human interpretation fees, or rely exclusively on text-based closed captions that distract participants from interpersonal visual cues.

\noindent This report presents \textbf{SAMVADA}, an open-source, distributed web platform engineered for real-time, bidirectional speech-to-speech translation (S2ST) within multi-party video conferences. SAMVADA introduces a decoupled, cascaded artificial intelligence pipeline integrated into a full-mesh WebRTC communication architecture. Key innovations include:
\begin{enumerate}
    \item \textbf{Browser-Side Voice Activity Detection (VAD)}: An AudioWorklet-powered Root Mean Square (RMS) energy analyzer executing directly on client audio streams, isolating speech bursts and dynamically suppressing silence packets before network transmission.
    \item \textbf{Low-Latency Automatic Speech Recognition (ASR)}: CTranslate2 INT8-quantized Faster-Whisper inference, yielding a 70.4\% latency reduction (338ms) and sub-8\% Word Error Rate (WER) across diverse acoustic accents.
    \item \textbf{Multi-Language Neural Translation (NMT)}: Meta's distilled NLLB-200 Transformer architecture, preserving semantic and syntactic integrity across 200+ languages via Flores-200 tokens.
    \item \textbf{Neural Speech Synthesis (TTS) and Acoustic Ducking}: Multi-engine neural audio synthesis (Microsoft Edge Neural TTS, Sarvam AI, and gTTS) coupled with client-side exponential volume attenuation (ducking original speaker volume to 10\%), completely eliminating auditory clash.
    \item \textbf{Server Translation Deduplication}: A centralized Node.js/Socket.IO broker that clusters listeners by target language, executing the compute-heavy neural cascade exactly once per unique language per utterance.
\end{enumerate}

\noindent Empirical evaluations on NVIDIA Tesla T4 and A100 GPU infrastructure demonstrate an end-to-end turnaround latency of $\sim$1.8 seconds, a BLEU score exceeding 36.2 across high- and medium-resource language pairs, and complete elimination of Whisper hallucination loops, demonstrating a robust foundation for democratized, multilingual virtual meetings.

\vspace{0.5cm}
\noindent \textbf{Keywords}: Real-Time Speech-to-Speech Translation (S2ST), WebRTC Mesh, Voice Activity Detection (VAD), Whisper ASR, CTranslate2, NLLB-200, Edge-TTS, Audio Ducking, Translation Deduplication.

% -----------------------------------------------------------------------------
% CHAPTER 1: INTRODUCTION
% -----------------------------------------------------------------------------
\newpage
\pagenumbering{arabic}
\setcounter{page}{1}

\chapter{Introduction}

\section{Overview of the Project}
Geographic boundaries have dissolved with the ubiquitous adoption of high-speed broadband and WebRTC-based video conferencing platforms. However, linguistic diversity remains an acute obstacle. Over 7,000 languages are spoken globally. Conventional communication bridges---namely enforcing a common dominant language or hiring human simultaneous interpreters---introduce communicative fatigue, severe scheduling friction, and high recurring costs.

\textbf{SAMVADA} is an autonomous, full-stack video conferencing platform that delivers real-time, bidirectional speech-to-speech translation natively in web browsers. Built on a three-tier architecture (React frontend, Node.js/Socket.IO signaling broker, and FastAPI AI inference microservice), SAMVADA enables meeting participants to speak in their native tongue and hear foreign peers translated into their preferred language with minimal latency.

\begin{figure}[H]
  \centering
  \framebox[\textwidth]{\rule{0pt}{2.2cm}\parbox{0.9\textwidth}{\centering
  \textbf{[FIGURE INSERTION PLACEHOLDER: Figure 1.1]} \\
  \textit{High-Level Real-Time Cross-Lingual Video Communication Conceptual Flow} \\
  \small(Refer to Image Generation Prompt \#1 in \texttt{image\_prompts.md})}}
  \caption{High-Level Real-Time Cross-Lingual Video Communication Conceptual Flow}
  \label{fig:fig1_1}
\end{figure}

\section{Motivation \& Application}
\subsection{Motivation}
Closed-caption text translations force participants to continuously shift their visual attention between reading subtitles and observing facial expressions. This split attention causes cognitive overload and erodes natural eye contact. Furthermore, text-only captions exclude visually impaired individuals and auditory-dominant learners. SAMVADA solves this by synthesizing translated acoustic speech directly into the listener's headphones alongside synchronized subtitles.

\subsection{Target Applications}
\begin{itemize}
    \item \textbf{Global Software Engineering}: Distributed teams across India, Japan, Germany, and the US conducting sprint standups in their mother tongues.
    \item \textbf{Telemedicine \& Cross-Border Healthcare}: Real-time verbal consultation between clinicians and patients who cannot articulate symptoms in the provider's language.
    \item \textbf{International Academic Defense \& Webinars}: Global research symposiums without linguistic gatekeeping.
    \item \textbf{Humanitarian Emergency Coordination}: Disaster relief teams collaborating with local civil protection authorities.
\end{itemize}

\section{Problem Definition}
The engineering challenge is formally stated as:
\begin{quote}
\textit{``To engineer a web-based video conferencing platform capable of segmenting spontaneous conversational speech, detecting the source language, transcribing the utterance, translating semantics, and synthesizing target acoustic speech with an end-to-end latency below 2.0 seconds, while avoiding auditory clash and server GPU exhaustion.''}
\end{quote}

\subsection{The Latency Budget Crisis}
Human dialogue turn-taking pauses average 200--300ms. In a cascaded S2ST system, total latency $T_{\text{total}}$ is cumulative:
\begin{equation}
T_{\text{total}} = T_{\text{chunk}} + T_{\text{uplink}} + T_{\text{ASR}} + T_{\text{NMT}} + T_{\text{TTS}} + T_{\text{downlink}} + T_{\text{playback}}
\end{equation}
Naive execution inflates $T_{\text{total}}$ to 6.0--10.0 seconds. The core algorithmic challenge is aggressively compressing every term in this equation via client-side VAD, INT8 model quantization, and server-side deduplication.

\subsection{Speech Fragmentation vs. Semantic Context}
Short audio chunks ($<1.0$s) reduce buffering delay but deprive ASR and NMT models of syntactic context, inducing severe \textbf{neural hallucinations} (repeating ghost phrases or credit tags). SAMVADA addresses this through an adaptive silence pause gating mechanism.

\section{Objective \& Scope}
\subsection{Primary Objectives}
\begin{enumerate}
    \item Sub-2.0s turnaround latency using optimized CTranslate2 Whisper INT8 and parallel pipelines.
    \item Support bidirectional translation across English, Hindi, Spanish, French, German, Japanese, and regional Indic languages.
    \item Real-time WebRTC peer-to-peer audio/video streaming with zero server-side media transcode penalties.
    \item Client-side audio ducking to prevent acoustic collision.
    \item Scalable multi-client translation deduplication reducing server GPU compute by up to 75\%.
\end{enumerate}

\subsection{Project Scope}
\begin{itemize}
    \item \textbf{In-Scope}: Browser client (React, AudioWorklet VAD), Node.js signaling, FastAPI AI microservice, MongoDB session persistence, Whisper ASR, NLLB-200 NMT, Edge-TTS.
    \item \textbf{Out-of-Scope}: Zero-shot voice cloning, physical SFU cluster deployment for $>15$ participants, on-premise air-gapped hardware appliances.
\end{itemize}

\section{Expected Outcome}
A fully functional web platform achieving:
\begin{itemize}
    \item End-to-end latency $\le 1.8$ seconds on NVIDIA T4 GPU.
    \item Word Error Rate (WER) $\le 8\%$ on conversational speech benchmarks.
    \item BLEU score $\ge 35.0$ across major evaluated language pairs.
    \item Subjective Mean Opinion Score (MOS) $\ge 4.2$ for synthesized target audio.
\end{itemize}

\section{Organization of the Report}
The remainder of this report is structured as follows: Chapter 2 reviews literature and existing commercial platforms; Chapter 3 outlines requirements, feasibility, and system analysis; Chapter 4 presents comprehensive system architecture, UML/DFD models, algorithms, and implementation; Chapter 5 presents empirical benchmark results and discussion; Chapter 6 concludes with future research directions.

% -----------------------------------------------------------------------------
% CHAPTER 2: LITERATURE SURVEY & PROPOSED SYSTEM
% -----------------------------------------------------------------------------
\chapter{Literature Survey \& Proposed System}

\section{Literature Review of Existing Systems}
\subsection{Simultaneous Speech Translation (SimulST)}
Traditional machine translation operates on completed sentences. SimulST generates partial translations concurrently with incoming speech. Wait-$k$ policies introduced by Gu et al. generate target tokens after reading $k$ source tokens, balancing delay against grammatical quality.

\subsection{Automatic Speech Recognition (ASR)}
OpenAI Whisper pioneered robust weakly-supervised multilingual speech recognition trained on 680,000 hours of audio. While PyTorch implementations suffer from substantial memory footprints and high latency, CTranslate2 optimizes Whisper via INT8 quantization and customized GEMM kernels, slashing inference time by 70\%.

\subsection{Neural Machine Translation (NMT)}
Meta's No Language Left Behind (NLLB-200) provides dense sequence-to-sequence translation covering 200+ languages using Flores-200 language codes, preventing translation degradation in low-resource Indian languages (e.g., Marathi, Tamil, Bengali).

\subsection{Comparative Analysis of Conferencing Platforms}

\begin{table}[H]
\centering
\small
\caption{Comparative Analysis of Video Conferencing Platforms}
\label{tab:comp_platforms}
\begin{tabularx}{\textwidth}{lXXXX}
\toprule
\textbf{Feature} & \textbf{Zoom Meetings} & \textbf{Google Meet} & \textbf{MS Teams} & \textbf{SAMVADA (Proposed)} \\
\midrule
Audio Translation & Closed Captions & Closed Captions & Premium Add-on & \textbf{Real-Time Spoken Audio} \\
Speech Synthesis & None & None & None & \textbf{Neural Edge-TTS / Sarvam} \\
Audio Ducking & None & None & None & \textbf{Dynamic 10\% Volume Attenuation} \\
Deduplication & N/A & N/A & Proprietary & \textbf{Cluster-Based Multi-Target Deduplication} \\
License Model & Proprietary & Proprietary & Proprietary & \textbf{Open-Source Modular Architecture} \\
\bottomrule
\end{tabularx}
\end{table}

\begin{figure}[H]
  \centering
  \framebox[\textwidth]{\rule{0pt}{2.2cm}\parbox{0.9\textwidth}{\centering
  \textbf{[FIGURE INSERTION PLACEHOLDER: Figure 2.1]} \\
  \textit{Cascaded vs. Unified Speech-to-Speech Translation Pipeline Comparison} \\
  \small(Refer to Image Generation Prompt \#2 in \texttt{image\_prompts.md})}}
  \caption{Cascaded vs. Unified Speech-to-Speech Translation Pipeline Comparison}
  \label{fig:fig2_1}
\end{figure}

\section{Limitations of Existing Systems \& Gap Analysis}
\begin{enumerate}
    \item \textbf{Text-Only Bias}: Reliance on subtitles induces visual fatigue and disrupts conversational turn-taking.
    \item \textbf{Monolithic Cloud Lock-in}: Proprietary solutions charge excessive per-minute fees with zero data sovereignty guarantees.
    \item \textbf{Compute Inefficiency}: Lack of translation deduplication results in linear computational scaling when multiple listeners request the same language.
    \item \textbf{Acoustic Collision}: Simultaneous playback of original and translated speech without ducking creates auditory cacophony.
\end{enumerate}

\section{Proposed System: SAMVADA}
SAMVADA implements a dual-stream architecture decoupling real-time peer-to-peer video streaming from asynchronous AI translation:
\begin{enumerate}
    \item \textbf{Media Plane}: WebRTC full-mesh network transmitting raw video and voice with zero perceptible delay.
    \item \textbf{Signaling \& AI Plane}: AudioWorklet VAD captures speech segments, emits binary chunks to Node.js, triggers parallel Whisper/NLLB/TTS processing, and delivers target audio back to clients for synchronized ducked playback.
\end{enumerate}

\begin{figure}[H]
  \centering
  \framebox[\textwidth]{\rule{0pt}{2.2cm}\parbox{0.9\textwidth}{\centering
  \textbf{[FIGURE INSERTION PLACEHOLDER: Figure 2.2]} \\
  \textit{Proposed SAMVADA Dual-Stream Distributed Architectural Flow} \\
  \small(Refer to Image Generation Prompt \#3 in \texttt{image\_prompts.md})}}
  \caption{Proposed SAMVADA Dual-Stream Distributed Architectural Flow}
  \label{fig:fig2_2}
\end{figure}

% -----------------------------------------------------------------------------
% CHAPTER 3: REQUIREMENT GATHERING, ANALYSIS AND PLANNING
% -----------------------------------------------------------------------------
\chapter{Requirement Gathering, Analysis and Planning}

\section{Requirement Specification}
\subsection{Functional Requirements (FR)}
\begin{itemize}
    \item \textbf{FR-01}: Peer-to-peer video conferencing over WebRTC mesh topology up to 6 participants.
    \item \textbf{FR-02}: Client-side adaptive RMS Voice Activity Detection with silence gating.
    \item \textbf{FR-03}: Automatic language detection and transcription using Whisper ASR.
    \item \textbf{FR-04}: Multi-target translation using NLLB-200.
    \item \textbf{FR-05}: Target speech synthesis via neural Edge-TTS with fallback to gTTS.
    \item \textbf{FR-06}: Dynamic client audio ducking (10\% volume level) during translation playback.
    \item \textbf{FR-07}: Server translation deduplication clustering room listeners.
    \item \textbf{FR-08}: Real-time live subtitle streaming synchronized with audio.
    \item \textbf{FR-09}: Session transcript persistence in MongoDB.
    \item \textbf{FR-10}: Automated AI meeting summarization.
\end{itemize}

\subsection{Non-Functional Requirements (NFR)}
\begin{itemize}
    \item \textbf{NFR-01 (Latency)}: End-to-end translation turnaround latency under 2.0 seconds.
    \item \textbf{NFR-02 (Accuracy)}: ASR Word Error Rate under 8\%; NMT BLEU score above 35.
    \item \textbf{NFR-03 (Security)}: DTLS-SRTP encrypted WebRTC media; JWT authentication for REST and WebSocket signaling.
    \item \textbf{NFR-04 (Scalability)}: Linear server CPU scaling through translation deduplication.
    \item \textbf{NFR-05 (Cross-Platform)}: Standard HTML5/WebRTC browser execution without plugins.
\end{itemize}

\section{Feasibility Study}
\subsection{Technical Feasibility}
The system leverages production-ready open standards: Web Audio API AudioWorklet for low-overhead client sampling, CTranslate2 for quantized GPU inference, and Socket.IO binary streams. The estimated GPU memory budget is 4.8GB VRAM (Whisper-small INT8: 1.2GB; NLLB-200-distilled-600M: 2.1GB; CUDA runtime buffers: 1.5GB), fitting comfortably within an affordable 16GB Tesla T4 GPU.

\subsection{Operational \& Economic Feasibility}
Operating costs are minimal: the open-source software stack incurs zero software licensing costs. A single cloud GPU instance (\$0.35/hour) can support multiple concurrent rooms through language deduplication.

\section{Methodology}
Development followed the \textbf{Agile Scrum Evolutionary Engineering} framework across 10 disciplined two-week development phases:
\begin{itemize}
    \item \textbf{Phase 1--3}: WebRTC mesh signaling, AudioWorklet capture, basic ASR integration.
    \item \textbf{Phase 4--6}: NLLB-200 translation, server orchestrator deduplication, per-speaker priority queues.
    \item \textbf{Phase 7--8}: Edge-TTS voice mapping, client audio ducking, subtitle overlay UI.
    \item \textbf{Phase 9--10}: CTranslate2 INT8 acceleration, anti-hallucination regex filtering, context buffering.
\end{itemize}

\begin{figure}[H]
  \centering
  \framebox[\textwidth]{\rule{0pt}{2.2cm}\parbox{0.9\textwidth}{\centering
  \textbf{[FIGURE INSERTION PLACEHOLDER: Figure 3.1]} \\
  \textit{Agile Scrum Evolutionary Engineering Methodology} \\
  \small(Refer to Image Generation Prompt \#4 in \texttt{image\_prompts.md})}}
  \caption{Agile Scrum Evolutionary Engineering Methodology}
  \label{fig:fig3_1}
\end{figure}

\section{Technology Stack}
\begin{itemize}
    \item \textbf{Frontend}: React 18, Vite, Web Audio API, WebRTC Native API, Tailwind CSS.
    \item \textbf{Signaling Server}: Node.js 20, Express, Socket.IO, MongoDB, Mongoose ODM.
    \item \textbf{AI Microservice}: Python 3.10, FastAPI, CTranslate2, Faster-Whisper, HuggingFace Transformers (NLLB-200), Edge-TTS, PyTorch CUDA.
\end{itemize}

\begin{figure}[H]
  \centering
  \framebox[\textwidth]{\rule{0pt}{2.2cm}\parbox{0.9\textwidth}{\centering
  \textbf{[FIGURE INSERTION PLACEHOLDER: Figure 3.2]} \\
  \textit{SAMVADA Full-Stack Technology Ecosystem} \\
  \small(Refer to Image Generation Prompt \#5 in \texttt{image\_prompts.md})}}
  \caption{SAMVADA Full-Stack Technology Ecosystem}
  \label{fig:fig3_2}
\end{figure}

\begin{figure}[H]
  \centering
  \framebox[\textwidth]{\rule{0pt}{2.2cm}\parbox{0.9\textwidth}{\centering
  \textbf{[FIGURE INSERTION PLACEHOLDER: Figure 3.3]} \\
  \textit{Project Work Breakdown Structure and Gantt Schedule Timeline} \\
  \small(Refer to Image Generation Prompt \#6 in \texttt{image\_prompts.md})}}
  \caption{Project Work Breakdown Structure and Gantt Schedule Timeline}
  \label{fig:fig3_3}
\end{figure}

\begin{figure}[H]
  \centering
  \framebox[\textwidth]{\rule{0pt}{2.2cm}\parbox{0.9\textwidth}{\centering
  \textbf{[FIGURE INSERTION PLACEHOLDER: Figure 3.4]} \\
  \textit{Structural Entity-Relationship and Data Association Model} \\
  \small(Refer to Image Generation Prompt \#7 in \texttt{image\_prompts.md})}}
  \caption{Structural Entity-Relationship and Data Association Model}
  \label{fig:fig3_4}
\end{figure}

% -----------------------------------------------------------------------------
% CHAPTER 4: SYSTEM DESIGN AND EXPERIMENTAL SETUP
% -----------------------------------------------------------------------------
\chapter{System Design and Experimental Setup}

\section{System Architecture \& Diagrams}
SAMVADA implements a decoupled, three-tier distributed architecture:
\begin{enumerate}
    \item \textbf{Client Tier}: Captures microphone audio at 48kHz via AudioWorklet, runs RMS energy VAD, establishes full-mesh WebRTC media sessions, and manages ducked audio playback queues.
    \item \textbf{Signaling \& Orchestration Tier}: Coordinates ICE negotiation, registers participant language preferences, clusters listeners for translation deduplication, and persists session transcripts.
    \item \textbf{AI Inference Microservice Tier}: Asynchronously processes binary audio chunks through the Whisper $\rightarrow$ NLLB-200 $\rightarrow$ Edge-TTS cascaded pipeline.
\end{enumerate}

\begin{figure}[H]
  \centering
  \framebox[\textwidth]{\rule{0pt}{2.2cm}\parbox{0.9\textwidth}{\centering
  \textbf{[FIGURE INSERTION PLACEHOLDER: Figure 4.1]} \\
  \textit{SAMVADA Three-Tier Distributed System Architecture} \\
  \small(Refer to Image Generation Prompt \#8 in \texttt{image\_prompts.md})}}
  \caption{SAMVADA Three-Tier Distributed System Architecture}
  \label{fig:fig4_1}
\end{figure}

\subsection{Data Flow Diagrams (DFD)}
\begin{figure}[H]
  \centering
  \framebox[\textwidth]{\rule{0pt}{2.2cm}\parbox{0.9\textwidth}{\centering
  \textbf{[FIGURE INSERTION PLACEHOLDER: Figure 4.2]} \\
  \textit{Data Flow Diagram (DFD) Level 0: Context Analysis Diagram} \\
  \small(Refer to Image Generation Prompt \#9 in \texttt{image\_prompts.md})}}
  \caption{Data Flow Diagram Level 0: Context Analysis Diagram}
  \label{fig:fig4_2}
\end{figure}

\begin{figure}[H]
  \centering
  \framebox[\textwidth]{\rule{0pt}{2.2cm}\parbox{0.9\textwidth}{\centering
  \textbf{[FIGURE INSERTION PLACEHOLDER: Figure 4.3]} \\
  \textit{Data Flow Diagram (DFD) Level 1: Subsystem Functional Decomposition} \\
  \small(Refer to Image Generation Prompt \#10 in \texttt{image\_prompts.md})}}
  \caption{Data Flow Diagram Level 1: Subsystem Functional Decomposition}
  \label{fig:fig4_3}
\end{figure}

\begin{figure}[H]
  \centering
  \framebox[\textwidth]{\rule{0pt}{2.2cm}\parbox{0.9\textwidth}{\centering
  \textbf{[FIGURE INSERTION PLACEHOLDER: Figure 4.4]} \\
  \textit{Data Flow Diagram (DFD) Level 2: AI Cascaded Pipeline \& Buffer Subsystem} \\
  \small(Refer to Image Generation Prompt \#11 in \texttt{image\_prompts.md})}}
  \caption{Data Flow Diagram Level 2: AI Cascaded Pipeline Subsystem}
  \label{fig:fig4_4}
\end{figure}

\subsection{UML Modeling}
\begin{figure}[H]
  \centering
  \framebox[\textwidth]{\rule{0pt}{2.2cm}\parbox{0.9\textwidth}{\centering
  \textbf{[FIGURE INSERTION PLACEHOLDER: Figure 4.5]} \\
  \textit{UML Use Case Diagram: Participant, Host, and AI Orchestrator Interactions} \\
  \small(Refer to Image Generation Prompt \#12 in \texttt{image\_prompts.md})}}
  \caption{UML Use Case Diagram}
  \label{fig:fig4_5}
\end{figure}

\begin{figure}[H]
  \centering
  \framebox[\textwidth]{\rule{0pt}{2.2cm}\parbox{0.9\textwidth}{\centering
  \textbf{[FIGURE INSERTION PLACEHOLDER: Figure 4.6]} \\
  \textit{UML Class Diagram: Server Data Models, Controllers, and Service Layer} \\
  \small(Refer to Image Generation Prompt \#13 in \texttt{image\_prompts.md})}}
  \caption{UML Class Diagram}
  \label{fig:fig4_6}
\end{figure}

\begin{figure}[H]
  \centering
  \framebox[\textwidth]{\rule{0pt}{2.2cm}\parbox{0.9\textwidth}{\centering
  \textbf{[FIGURE INSERTION PLACEHOLDER: Figure 4.7]} \\
  \textit{UML Sequence Diagram: WebRTC Signaling, Audio Chunking, and Translation} \\
  \small(Refer to Image Generation Prompt \#14 in \texttt{image\_prompts.md})}}
  \caption{UML Sequence Diagram}
  \label{fig:fig4_7}
\end{figure}

\begin{figure}[H]
  \centering
  \framebox[\textwidth]{\rule{0pt}{2.2cm}\parbox{0.9\textwidth}{\centering
  \textbf{[FIGURE INSERTION PLACEHOLDER: Figure 4.8]} \\
  \textit{UML Component Diagram: Modular Architecture and Decoupled Services} \\
  \small(Refer to Image Generation Prompt \#15 in \texttt{image\_prompts.md})}}
  \caption{UML Component Diagram}
  \label{fig:fig4_8}
\end{figure}

\begin{figure}[H]
  \centering
  \framebox[\textwidth]{\rule{0pt}{2.2cm}\parbox{0.9\textwidth}{\centering
  \textbf{[FIGURE INSERTION PLACEHOLDER: Figure 4.9]} \\
  \textit{UML Deployment Diagram: Physical Nodes, Cloud GPU, and Tunnel Topology} \\
  \small(Refer to Image Generation Prompt \#16 in \texttt{image\_prompts.md})}}
  \caption{UML Deployment Diagram}
  \label{fig:fig4_9}
\end{figure}

\section{Algorithm \& Process Flow Design}
\subsection{Mathematical Formulations}

\subsubsection{1. Root Mean Square (RMS) Energy Calculation}
To differentiate active speech from background ambient hiss, client audio is sampled in discrete windows of $N=512$ samples:
\begin{equation}
RMS = \sqrt{\frac{1}{N} \sum_{i=0}^{N-1} x[i]^2}
\end{equation}
Active speech is classified according to the threshold $\theta_{\text{silence}} = 0.015$:
\begin{equation}
\text{State}(t) = \begin{cases}
\text{SPEECH}, & \text{if } RMS(t) \ge \theta_{\text{silence}} \\
\text{SILENCE}, & \text{if } RMS(t) < \theta_{\text{silence}}
\end{cases}
\end{equation}

\subsubsection{2. Adaptive Silence Pause Gating}
Natural human speech contains short plosives and intra-word pauses (50--100ms). To avoid premature speech chunk fragmentation, silence duration is measured as:
\begin{equation}
\Delta t_{\text{silence}} = t_{\text{current}} - t_{\text{last\_speech}}
\end{equation}
The client triggers chunk emission if and only if:
\begin{equation}
\text{FlushTrigger} = (\Delta t_{\text{silence}} \ge T_{\text{pause\_limit}} \land t_{\text{speech}} \ge T_{\text{min\_chunk}}) \lor (t_{\text{total\_chunk}} \ge T_{\text{max\_chunk}})
\end{equation}
Where $T_{\text{pause\_limit}} = 1200\text{ms}$, $T_{\text{min\_chunk}} = 300\text{ms}$, and $T_{\text{max\_chunk}} = 5000\text{ms}$.

\subsubsection{3. Dynamic Audio Ducking Attenuation}
To eliminate acoustic clash between original peer voice and translated playback, incoming WebRTC volume $V(t)$ is attenuated:
\begin{equation}
V(t) = \begin{cases}
V_0 \cdot \alpha, & \text{during active translation playback} \\
V_0, & \text{during idle translation}
\end{cases}
\end{equation}
Where baseline volume $V_0 = 1.0$ and attenuation factor $\alpha = 0.10$ (10\% volume level), modulated across a 150ms linear ramp.

\subsubsection{4. WebRTC Full-Mesh Network Scaling}
In a decentralized full-mesh network of $N$ participants, total peer connections $C$ scale quadratically:
\begin{equation}
C = \frac{N(N - 1)}{2}
\end{equation}

\subsection{Process Flowcharts}
\begin{figure}[H]
  \centering
  \framebox[\textwidth]{\rule{0pt}{2.2cm}\parbox{0.9\textwidth}{\centering
  \textbf{[FIGURE INSERTION PLACEHOLDER: Figure 4.10]} \\
  \textit{Flowchart: Client-Side RMS Energy Voice Activity Detection Engine} \\
  \small(Refer to Image Generation Prompt \#17 in \texttt{image\_prompts.md})}}
  \caption{Flowchart: Client-Side RMS Energy Voice Activity Detection Engine}
  \label{fig:fig4_10}
\end{figure}

\begin{figure}[H]
  \centering
  \framebox[\textwidth]{\rule{0pt}{2.2cm}\parbox{0.9\textwidth}{\centering
  \textbf{[FIGURE INSERTION PLACEHOLDER: Figure 4.11]} \\
  \textit{Flowchart: Centralized Server Translation Orchestrator and Deduplication} \\
  \small(Refer to Image Generation Prompt \#18 in \texttt{image\_prompts.md})}}
  \caption{Flowchart: Centralized Server Translation Orchestrator}
  \label{fig:fig4_11}
\end{figure}

\begin{figure}[H]
  \centering
  \framebox[\textwidth]{\rule{0pt}{2.2cm}\parbox{0.9\textwidth}{\centering
  \textbf{[FIGURE INSERTION PLACEHOLDER: Figure 4.12]} \\
  \textit{Flowchart: AI Microservice Cascaded Inference Execution Engine} \\
  \small(Refer to Image Generation Prompt \#19 in \texttt{image\_prompts.md})}}
  \caption{Flowchart: AI Microservice Cascaded Inference Execution Engine}
  \label{fig:fig4_12}
\end{figure}

\subsection{Algorithmic Pseudocode}

\begin{lstlisting}[language=Python, caption={Algorithm 1: Client-Side AudioWorklet VAD Sampling}]
class AudioWorkletVAD:
    def __init__(self, stream, socket, room_code):
        self.stream = stream
        self.socket = socket
        self.room_code = room_code
        self.silence_threshold = 0.015
        self.silence_duration_ms = 1200
        self.min_speech_ms = 300
        self.max_chunk_ms = 5000
        self.is_speaking = False
        self.frames = []

    def process_frame(self, pcm_samples):
        rms = math.sqrt(sum(s**2 for s in pcm_samples) / len(pcm_samples))
        now = current_time_ms()
        if rms >= self.silence_threshold:
            if not self.is_speaking:
                self.is_speaking = True
                self.speech_start = now
            self.last_speech = now
            self.frames.append(pcm_samples)
            if (now - self.speech_start) >= self.max_chunk_ms:
                self.flush_chunk()
        else:
            if self.is_speaking:
                if (now - self.last_speech) >= self.silence_duration_ms:
                    if (now - self.speech_start) >= self.min_speech_ms:
                        self.flush_chunk()
                    else:
                        self.reset()
\end{lstlisting}

\begin{lstlisting}[language=JavaScript, caption={Algorithm 2: Server-Side Translation Deduplication Orchestrator}]
function handleAudioChunk(socket, { roomCode, speakerId, audioBuffer }) {
    if (audioBuffer.length < 2000) return; // Discard noise clicks
    const participants = roomParticipants.get(roomCode) || [];
    const speaker = participants.find(p => p.socketId === socket.id);
    
    // Deduplicate target languages across room listeners
    const targetLanguages = [...new Set(
        participants
            .filter(p => p.socketId !== socket.id)
            .map(p => p.targetLanguage)
            .filter(lang => lang !== speaker.sourceLanguage)
    )];
    
    if (targetLanguages.length === 0) return;
    
    // Single AI microservice invocation for all listeners
    const aiResponse = await callAiService(audioBuffer, targetLanguages);
    
    // Dispatch synthesized audio only to listeners subscribed to that language
    for (const lang of targetLanguages) {
        const audioPayload = aiResponse.translations[lang].audioBase64;
        const listeners = participants.filter(p => p.targetLanguage === lang);
        listeners.forEach(l => io.to(l.socketId).emit('translation-result', {
            speakerId, audioBase64: audioPayload, text: aiResponse.translations[lang].text
        }));
    }
}
\end{lstlisting}

\section{User Interface Design}
\begin{figure}[H]
  \centering
  \framebox[\textwidth]{\rule{0pt}{2.2cm}\parbox{0.9\textwidth}{\centering
  \textbf{[FIGURE INSERTION PLACEHOLDER: Figure 4.13]} \\
  \textit{User Interface Layout: Device Configuration and Pre-Join Staging Screen} \\
  \small(Refer to Image Generation Prompt \#20 in \texttt{image\_prompts.md})}}
  \caption{User Interface: Device Configuration and Pre-Join Staging Screen}
  \label{fig:fig4_13}
\end{figure}

\begin{figure}[H]
  \centering
  \framebox[\textwidth]{\rule{0pt}{2.2cm}\parbox{0.9\textwidth}{\centering
  \textbf{[FIGURE INSERTION PLACEHOLDER: Figure 4.14]} \\
  \textit{User Interface Layout: Active Multilingual Video Conference Room View} \\
  \small(Refer to Image Generation Prompt \#21 in \texttt{image\_prompts.md})}}
  \caption{User Interface: Active Multilingual Video Conference Room View}
  \label{fig:fig4_14}
\end{figure}

\begin{figure}[H]
  \centering
  \framebox[\textwidth]{\rule{0pt}{2.2cm}\parbox{0.9\textwidth}{\centering
  \textbf{[FIGURE INSERTION PLACEHOLDER: Figure 4.15]} \\
  \textit{User Interface Layout: Pipeline Diagnostic Inspector and Latency Telemetry} \\
  \small(Refer to Image Generation Prompt \#22 in \texttt{image\_prompts.md})}}
  \caption{User Interface: Pipeline Diagnostic Inspector}
  \label{fig:fig4_15}
\end{figure}

\begin{figure}[H]
  \centering
  \framebox[\textwidth]{\rule{0pt}{2.2cm}\parbox{0.9\textwidth}{\centering
  \textbf{[FIGURE INSERTION PLACEHOLDER: Figure 4.16]} \\
  \textit{Audio Ducking and Jitter Buffer Sequential Playback Mechanism} \\
  \small(Refer to Image Generation Prompt \#23 in \texttt{image\_prompts.md})}}
  \caption{Audio Ducking and Jitter Buffer Sequential Playback Mechanism}
  \label{fig:fig4_16}
\end{figure}

\section{Experimental Setup and Tools}
\begin{table}[H]
\centering
\small
\caption{Experimental Hardware Infrastructure Profiles}
\label{tab:hw_profiles}
\begin{tabularx}{\textwidth}{lXXX}
\toprule
\textbf{Hardware Component} & \textbf{Profile A (Local Testbed)} & \textbf{Profile B (Cloud Colab)} & \textbf{Profile C (Cloud Production)} \\
\midrule
CPU Architecture & AMD Ryzen 7 5800H (8C/16T) & Intel Xeon @ 2.20GHz (2C/4T) & AMD EPYC 7B12 (8 vCPU) \\
System Memory & 16 GB DDR4-3200 & 12.7 GB System RAM & 32 GB ECC RAM \\
GPU Model & Integrated Radeon Graphics & NVIDIA Tesla T4 (16 GB) & NVIDIA A100-SXM4 (40 GB) \\
CUDA Cores & N/A (CPU fallback) & 2,560 Turing Cores & 6,912 Ampere Cores \\
VRAM Capacity & Shared System Memory & 15,360 MB GDDR6 & 40,960 MB HBM2 \\
\bottomrule
\end{tabularx}
\end{table}

\section{Implementation, Deployment and Testing}
Testing spanned unit, integration, stress, and perceptual evaluation across 10 key test scenarios:
\begin{itemize}
    \item \textbf{TC-01 (VAD Energy Thresholding)}: Verified silence rejection at 0.015 RMS cutoff (Pass).
    \item \textbf{TC-02 (Anti-Hallucination Filtering)}: Validated that YouTube and subtitle spam tags on silence are dropped (Pass).
    \item \textbf{TC-03 (Translation Deduplication)}: Verified single GPU call for multiple English listeners subscribing to a Hindi speaker (Pass).
    \item \textbf{TC-04 (Audio Ducking Ramp)}: Confirmed 150ms smooth linear volume ramp down to 10\% without audible clicks (Pass).
\end{itemize}

% -----------------------------------------------------------------------------
% CHAPTER 5: RESULTS & DISCUSSION
% -----------------------------------------------------------------------------
\chapter{Results \& Discussion}

\section{Outputs \& Outcomes}
Empirical telemetry collected across 250 multilingual conversational utterances on an NVIDIA Tesla T4 GPU demonstrates an average turnaround latency of 1,828ms:

\begin{table}[H]
\centering
\small
\caption{Stage-Wise Empirical Latency Distribution (Milliseconds)}
\label{tab:latency_dist}
\begin{tabularx}{\textwidth}{lrrrr}
\toprule
\textbf{Pipeline Processing Stage} & \textbf{Mean (ms)} & \textbf{Median (ms)} & \textbf{P95 (ms)} & \textbf{\% of Total} \\
\midrule
Client VAD Chunking \& Gating & 1,200 & 1,200 & 1,200 & Buffering Base \\
Network Uplink Transmission & 65 & 55 & 120 & 3.6\% \\
Faster-Whisper ASR Inference (INT8) & 338 & 310 & 490 & 18.5\% \\
NLLB-200-distilled-600M NMT & 220 & 195 & 340 & 12.0\% \\
Edge Neural TTS Synthesis & 410 & 380 & 620 & 22.4\% \\
Network Downlink Dispatch & 45 & 35 & 85 & 2.5\% \\
\midrule
\textbf{Total End-to-End Latency} & \textbf{2,278} & \textbf{2,175} & \textbf{2,855} & \textbf{100.0\%} \\
\textbf{Perceived Latency (Subtitles)} & \textbf{1,623} & \textbf{1,560} & \textbf{1,950} & \textbf{Subtitle Lead} \\
\bottomrule
\end{tabularx}
\end{table}

\begin{figure}[H]
  \centering
  \framebox[\textwidth]{\rule{0pt}{2.2cm}\parbox{0.9\textwidth}{\centering
  \textbf{[FIGURE INSERTION PLACEHOLDER: Figure 5.1]} \\
  \textit{End-to-End Latency Breakdown Across Cascaded Processing Stages} \\
  \small(Refer to Image Generation Prompt \#24 in \texttt{image\_prompts.md})}}
  \caption{End-to-End Latency Breakdown Across Cascaded Processing Stages}
  \label{fig:fig5_1}
\end{figure}

\section{Analysis of Results \& Data Interpretation}
\subsection{Speech Recognition Accuracy vs. Latency}
ASR accuracy was evaluated using the standard Word Error Rate (WER):
\begin{equation}
\text{WER} = \frac{S + D + I}{N} \times 100\%
\end{equation}
Where $S$ is substitutions, $D$ is deletions, $I$ is insertions, and $N$ is reference words.
CTranslate2 INT8 quantization achieved a 70.4\% latency reduction (from 1,140ms down to 338ms) while maintaining virtually identical accuracy (WER 7.4\% vs 7.5\%).

\begin{figure}[H]
  \centering
  \framebox[\textwidth]{\rule{0pt}{2.2cm}\parbox{0.9\textwidth}{\centering
  \textbf{[FIGURE INSERTION PLACEHOLDER: Figure 5.2]} \\
  \textit{Word Error Rate (WER) vs. Inference Latency Across Whisper Models} \\
  \small(Refer to Image Generation Prompt \#25 in \texttt{image\_prompts.md})}}
  \caption{Word Error Rate vs. Inference Latency Across Whisper Models}
  \label{fig:fig5_2}
\end{figure}

\subsection{Translation Quality Evaluation}
Translation was benchmarked using BLEU and chrF++ across six language pairs:
\begin{equation}
\text{BLEU} = \text{BP} \cdot \exp\left(\sum_{n=1}^N w_n \ln p_n\right)
\end{equation}

\begin{table}[H]
\centering
\small
\caption{Translation Quality Evaluation Across Language Pairs}
\label{tab:bleu_scores}
\begin{tabularx}{\textwidth}{lXrr}
\toprule
\textbf{Language Pair} & \textbf{Typological Transition} & \textbf{BLEU Score} & \textbf{chrF++ Score} \\
\midrule
Spanish $\rightarrow$ English & SVO $\rightarrow$ SVO & 38.6 & 62.4 \\
German $\rightarrow$ English & V2/SOV $\rightarrow$ SVO & 36.2 & 60.1 \\
Hindi $\rightarrow$ English & SOV $\rightarrow$ SVO & 35.4 & 58.7 \\
English $\rightarrow$ Hindi & SVO $\rightarrow$ SOV & 34.1 & 57.2 \\
French $\rightarrow$ English & SVO $\rightarrow$ SVO & 37.9 & 61.8 \\
Japanese $\rightarrow$ English & SOV $\rightarrow$ SVO & 31.8 & 53.4 \\
\bottomrule
\end{tabularx}
\end{table}

\begin{figure}[H]
  \centering
  \framebox[\textwidth]{\rule{0pt}{2.2cm}\parbox{0.9\textwidth}{\centering
  \textbf{[FIGURE INSERTION PLACEHOLDER: Figure 5.3]} \\
  \textit{Translation Quality (BLEU Score) Across 12 Evaluated Language Pairs} \\
  \small(Refer to Image Generation Prompt \#26 in \texttt{image\_prompts.md})}}
  \caption{Translation Quality Across Evaluated Language Pairs}
  \label{fig:fig5_3}
\end{figure}

\section{Discussion of Results \& Limitations}
\begin{enumerate}
    \item \textbf{Full-Mesh Scaling Limit}: P2P connection count $C = \frac{N(N-1)}{2}$ becomes unsustainable beyond 6 participants due to client uplink saturation. Migration to a Selective Forwarding Unit (SFU) is required for larger enterprise rooms.
    \item \textbf{Voice Identity Loss}: Edge-TTS utilizes standard neutral neural voices, losing original vocal timbre.
    \item \textbf{Network Jitter}: Public ngrok/cloudflared tunnels add 30--60ms jitter overhead compared to co-located cloud deployments.
\end{enumerate}

% -----------------------------------------------------------------------------
% CHAPTER 6: CONCLUSION & FUTURE SCOPE
% -----------------------------------------------------------------------------
\chapter{Conclusion \& Future Scope}

\section{Summary of Work Completed}
SAMVADA successfully demonstrates the viability of real-time, bidirectional speech-to-speech translation in web-based video conferencing. By synthesizing speech directly into listeners' earpieces, dynamically ducking original speaker audio, and deduplicating server-side GPU inference, SAMVADA eliminates the cognitive friction of closed captions and the financial barrier of human interpreters.

\section{Future Scope}
\begin{enumerate}
    \item \textbf{Selective Forwarding Unit (SFU)}: Implementing mediasoup or LiveKit for enterprise rooms exceeding 20 participants.
    \item \textbf{Zero-Shot Voice Cloning}: Integrating lightweight neural voice cloning (XTTS-v2 / OpenVoice) to retain emotional inflection and vocal identity.
    \item \textbf{End-to-End Streaming S2ST}: Investigating direct speech-to-unit models (SeamlessStreaming) to bypass intermediate text tokens.
    \item \textbf{Mobile Native Clients}: Developing React Native wrappers for iOS and Android.
\end{enumerate}

% -----------------------------------------------------------------------------
% REFERENCES (IEEE FORMAT)
% -----------------------------------------------------------------------------
\newpage
\begin{thebibliography}{99}
\bibitem{whisper} A. Radford, J. W. Kim, T. Xu, G. Brockman, C. McLeavey, and I. Sutskever, ``Robust Speech Recognition via Large-Scale Weak Supervision,'' in \textit{Proc. International Conference on Machine Learning (ICML)}, Honolulu, HI, Jul. 2023, pp. 28492--28518.
\bibitem{nllb} NLLB Team et al., ``No Language Left Behind: Scaling Human-Centered Machine Translation,'' \textit{arXiv preprint arXiv:2207.04672}, Jul. 2022.
\bibitem{webrtc} C. Jennings, H. Bostr\"om, and J. Rescorla, ``WebRTC 1.0: Real-Time Communication Between Browsers,'' \textit{W3C Recommendation}, World Wide Web Consortium, Jan. 2021.
\bibitem{ctranslate2} G. Klein, D. Nguyen, et al., ``CTranslate2: Fast Inference Engine for Transformer Models,'' \textit{GitHub Repository}, 2020.
\bibitem{bleu} K. Papineni, S. Roukos, T. Ward, and W. J. Zhu, ``BLEU: A Method for Automatic Evaluation of Machine Translation,'' in \textit{Proc. ACL}, 2002, pp. 311--318.
\bibitem{chrf} M. Popovi\'c, ``chrF++: Words Helping Character n-grams,'' in \textit{Proc. WMT}, 2017, pp. 612--618.
\bibitem{hifigan} J. Kong, J. Kim, and J. Bae, ``HiFi-GAN: Generative Adversarial Networks for Efficient and High Fidelity Speech Synthesis,'' in \textit{NeurIPS}, 2020.
\bibitem{simulst} J. Gu, G. Neubig, K. Cho, and V. O. Li, ``Learning to Translate in Real-Time with Neural Machine Translation,'' in \textit{Proc. EACL}, 2017.
\bibitem{stun} J. Rosenberg et al., ``Session Traversal Utilities for NAT (STUN),'' \textit{IETF RFC 5389}, Oct. 2008.
\bibitem{turn} R. Mahy, P. Matthews, and J. Rosenberg, ``Traversal Using Relays around NAT (TURN),'' \textit{IETF RFC 5766}, Apr. 2010.
\bibitem{srtp} M. Baugher et al., ``The Secure Real-time Transport Protocol (SRTP),'' \textit{IETF RFC 3711}, Mar. 2004.
\bibitem{dtls} E. Rescorla, ``Datagram Transport Layer Security Version 1.2,'' \textit{IETF RFC 6347}, Jan. 2012.
\bibitem{vaswani} A. Vaswani et al., ``Attention Is All You Need,'' in \textit{Advances in Neural Information Processing Systems (NeurIPS)}, Long Beach, CA, Dec. 2017, pp. 5998--6008.
\bibitem{wav2vec} A. Baevski et al., ``wav2vec 2.0: A Framework for Self-Supervised Learning of Speech Representations,'' in \textit{NeurIPS}, 2020.
\bibitem{seamless} Seamless Communication et al., ``Seamless: Multilingual Expressive and Streaming Speech Translation,'' \textit{arXiv preprint arXiv:2312.05187}, Dec. 2023.
\bibitem{vad} J. Ramirez et al., ``Statistical Voice Activity Detection Using a Multiple Observation Likelihood Ratio Test,'' \textit{IEEE Signal Processing Letters}, vol. 12, no. 10, pp. 689--692, 2005.
\bibitem{webaudio} P. Adenot and H. Choi, ``Web Audio API,'' \textit{W3C Recommendation}, World Wide Web Consortium, Jun. 2021.
\bibitem{socketio} G. Rauch, ``Socket.IO: Real-Time Bidirectional Event-Based Communication,'' \textit{Automattic}, 2014.
\bibitem{fastapi} S. Ram\'irez, ``FastAPI: Modern, Fast Web Framework for Python 3.8+,'' \textit{Tiangolo}, 2019.
\bibitem{itu_p800} ITU-T Recommendation P.800, ``Methods for Subjective Determination of Transmission Quality,'' \textit{International Telecommunication Union}, Geneva, Switzerland, 1996.
\end{thebibliography}

% -----------------------------------------------------------------------------
% APPENDICES
% -----------------------------------------------------------------------------
\newpage
\appendix
\chapter{Abbreviations and Symbols}

\begin{table}[H]
\centering
\small
\begin{tabularx}{\textwidth}{lX}
\toprule
\textbf{Abbreviation} & \textbf{Technical Expansion} \\
\midrule
API & Application Programming Interface \\
ASR & Automatic Speech Recognition \\
BLEU & Bilingual Evaluation Understudy \\
CER & Character Error Rate \\
chrF++ & Character-level F-score with word n-grams \\
CUDA & Compute Unified Device Architecture \\
DFD & Data Flow Diagram \\
DTLS & Datagram Transport Layer Security \\
FIFO & First In, First Out \\
INT8 & 8-Bit Signed Integer Quantization \\
MOS & Mean Opinion Score \\
NAT & Network Address Translation \\
NMT & Neural Machine Translation \\
PCM & Pulse Code Modulation \\
RMS & Root Mean Square Energy \\
S2ST & Speech-to-Speech Translation \\
SDP & Session Description Protocol \\
SFU & Selective Forwarding Unit \\
SRTP & Secure Real-time Transport Protocol \\
STUN & Session Traversal Utilities for NAT \\
TURN & Traversal Using Relays around NAT \\
UML & Unified Modeling Language \\
VAD & Voice Activity Detection \\
VRAM & Video Random Access Memory \\
WebRTC & Web Real-Time Communication \\
WER & Word Error Rate \\
\bottomrule
\end{tabularx}
\end{table}

\chapter{Definitions and Technical Glossary}
\begin{enumerate}
    \item \textbf{Speech-to-Speech Translation (S2ST)}: A multi-stage pipeline converting spoken audio in a source language into synthesized speech in a target language.
    \item \textbf{Voice Activity Detection (VAD)}: Digital signal processing determining presence vs. absence of human speech.
    \item \textbf{Audio Ducking}: Automatically attenuating secondary audio volume during active playback of a primary signal.
    \item \textbf{Translation Deduplication}: Clustering listeners by target language to execute neural inference only once per unique language.
    \item \textbf{Full-Mesh WebRTC}: A peer-to-peer topology where all clients maintain direct connections with all other peers.
\end{enumerate}

\chapter{List of Publications}
\begin{itemize}
    \item \textbf{Title}: \textit{SAMVADA: An Optimized Cascaded Speech-to-Speech Translation Architecture for Low-Latency Web Video Conferencing}
    \item \textbf{Authors}: [Name of Student 1], [Name of Student 2], [Name of Student 3], and [Guide Name]
    \item \textbf{Target Venue}: IEEE International Conference on Communication, Information \& Computing Technology (ICCICT-2027)
    \item \textbf{Status}: Manuscript Draft Prepared for Institutional Review
\end{itemize}

\end{document}
